\section{Confined renewal rewards}
\label{R:alt:section}

The length tail of one irreducible does not specify where its expected
length is accumulated. We now retain that length in two kinds of spatial
region. The polynomial vacuum estimates give an irreducible reward in a
logarithmically enlarged box and an expected number of completed vertices
before exit. The disk representation gives the expected completed-prefix
mass inside a ball whose radius is a fixed multiple of the height scale.
These conclusions use the common renewal identities and their own finite
inputs, stated below; the weaker hypotheses in the first argument do not
require the sharp length deficit of Section~\ref{R:sec:renewal}.

\subsection{Length rewards with logarithmic confinement}
\label{R:alt:pv}

We use the honeycomb lattice dual to a triangular tiling of side one.
Successive horizontal cuts are at Euclidean distance $\sqrt3/2$; their
integer separation is the height.  As before, a port has no weight, and a
path visiting $L$ vertices has weight $\rho^L$.  Write $b_h$ for the critical
mass of bridges from a fixed bottom port to all top ports at height $h$,
$b_0=1$, and $m_h$ for the same mass with an insertion of $L$, $m_0=0$.
The polynomial vacuum representation gives the following finite inputs
\cite[Theorem~1.1, Proposition~1.3 and Lemma~8.2]{RenewalVacuumCompanion}:
\begin{equation}\label{R:alt:pv:inputs}
 b_h=h^{-1/4+o(1)},\qquad m_h=h^{13/12+o(1)}.
\end{equation}
For a fixed-height bridge let $S$ be its largest absolute horizontal
displacement from its source.  There are absolute positive $c,C$ such that
\begin{equation}\label{R:alt:pv:lateral}
 \sum_{\gamma:\,H(\gamma)=h,\ S(\gamma)>t}\rho^{L(\gamma)}
 \le C e^{-ct/(h+1)}.
\end{equation}
The same bound holds on the subclass of irreducibles.  The mixed exit law in
the strip has total first length mass at most $h^{13/12+o(1)}$; the mass of
half-plane arches not contained in that strip is $O(b_h)$.
For the latter assertion, let $A_h$ be the strip return-arch mass and
$c_0=\cos(3\pi/8)$.  The strip flux identity in
\cite[Section~2]{RenewalVacuumCompanion} is $b_h+c_0A_h=1$.
Since $b_h\to0$, monotone exhaustion gives $c_0A_\infty=1$ and
$A_\infty-A_h=b_h/c_0$.  All these are fixed-source, unnormalized masses.  The logarithmic confinement in
\eqref{R:alt:pv:lateral} will be used explicitly below.

Let $p$ be the irreducible probability from
Proposition~\ref{R:renew:factorization}.  Its normalization proof only
uses the concatenation identity and divergence of $\sum_hb_h$, so
\eqref{R:alt:pv:inputs} suffices for this application.  Under $p$, write
$H,L,D$ for height,
vertex length and diameter, and $S$ for largest horizontal displacement.
Put $\beta=3/4$ and $d=4/3$.  The following reward estimate is also useful
when the terminal port is subsequently prescribed.

\begin{proposition}[Truncated and discounted irreducible rewards]
\label{R:alt:pv:reward}
For this irreducible law,
\begin{align}
 p(H>r)&=r^{-\beta+o(1)},&
 \mathbb E_p[L;H\le r]&=r^{d-\beta+o(1)},
 \label{R:alt:pv:truncated}\\
 \mathbb E_p[Le^{-H/r}]&=r^{d-\beta+o(1)}.
 \label{R:alt:pv:discounted}
\end{align}
The lower bound in \eqref{R:alt:pv:discounted} is retained after imposing
$H\le r\log^2r$ and $S\le r\log^4r$.  Furthermore,
\begin{equation}\label{R:alt:pv:tails}
 p(L>n)\le n^{-9/16+o(1)},\qquad
 p(D>r)\le r^{-3/4+o(1)},\qquad
 1-\mathbb E_p e^{-L/n}\ge n^{-17/24-o(1)}.
\end{equation}
\end{proposition}
\begin{proof}
For $s>0$, the renewal and reward identities give
\[
 B(s):=\sum_{h\ge0}b_he^{-sh}=\frac1{1-\mathbb E_p e^{-sH}},\qquad
 \sum_{h\ge0}m_he^{-sh}=B(s)^2\mathbb E_p[Le^{-sH}].
\]
Positive coefficient summation of \eqref{R:alt:pv:inputs} yields
$B(1/r)=r^{\beta+o(1)}$ and
$\sum m_he^{-h/r}=r^{d+\beta+o(1)}$.  This proves
\eqref{R:alt:pv:discounted}.  The height deficit is
$r^{-\beta+o(1)}$, so $p(H>r)\le r^{-\beta+o(1)}$ and
$\mathbb E[H;H\le r]\le r^{1-\beta+o(1)}$.
To obtain the reverse height-tail bound, fix $\delta>0$ and use the deficit
at scale $r^{1+\delta}$.  Heights at most $r$ contribute at most
$r^{-\beta-\delta+o(1)}$, whereas that deficit is
$r^{-\beta(1+\delta)+o(1)}$.  The former is negligible because $\beta<1$.
Thus $p(H>r)\ge r^{-\beta(1+\delta)-o(1)}$; let $\delta$ decrease to zero.

Write $j_h=\mathbb E_p[L;H=h]$.  Coefficientwise the reward identity is
$m=b*j*b$.  At height $3r$, the convolution multiplier of each $j_h$ with
$h\le r$ is
\[
 (b*b)_{3r-h}\ge r^{2\beta-1-o(1)},
\]
because order $r$ terms have both indices comparable to $r$.
Since $m_{3r}\le r^{d+\beta-1+o(1)}$, this proves the upper bound in
\eqref{R:alt:pv:truncated}.  Splitting $p(L>n)$ according as
$H>n^{1/d}$ or not gives the first bound in \eqref{R:alt:pv:tails}.

We verify that the discounted reward is spatially retained.  The elementary
packing bound for a self-avoiding path in a strip is
$L\le C h(1+S)$ at height $h\ge1$.
Integrating \eqref{R:alt:pv:lateral} therefore gives, after decreasing $c$,
\[
 \mathbb E_p[L;H=h,S>t]
 \le C h(1+t+h)e^{-ct/(h+1)}.
\]
For $h>r\log^2r$, use $j_h\le m_h$ and the factor $e^{-h/r}$;
the resulting tail is smaller than every inverse power of $r$.
For $h\le r\log^2r$ and $S>r\log^4r$, the last display has the same
property after summing over $h$.  Thus removing these two tails retains
$r^{d-\beta+o(1)}$.  Evaluating this conclusion at $r=h^{1-\delta}$
gives $\sum_{j\le h}j_j\ge h^{d-\beta-o(1)}$ after letting
$\delta\downarrow0$.  This proves the other half of
\eqref{R:alt:pv:truncated}.

On the retained set $L=O(r^2\log^6r)$.  Take
$r=n^{1/2-\delta}$, so eventually $L\le n$ on that set.  Since
$1-e^{-L/n}\ge(1-e^{-1})L/n$ for $L\le n$, its retained reward gives
\[
 1-\mathbb E_p e^{-L/n}
 \ge n^{-1}r^{7/12-o(1)}
 =n^{-17/24-(7/12)\delta-o(1)}.
\]
Let $\delta\downarrow0$.  Finally, split the event $D>r$ at
$H=r/\log^2r$.  The height tail costs $r^{-\beta+o(1)}$; on the
complement, diameter greater than $r$ forces $S\ge c_0r$.
Summing \eqref{R:alt:pv:lateral} over the remaining heights is
superpolynomially small.  This proves the diameter bound.
\end{proof}

To distinguish a completed prefix from the part retained before an exit,
write $\eta_k$ for the $k$th irreducible, $H_k=H(\eta_k)$, $L_k=L(\eta_k)$,
$S_H(k)=\sum_{i\le k}H_i$, and $\Gamma_k=\eta_1\circ\cdots\circ\eta_k$.
For $h,w>0$ put
\[
 V_h=\sum_{k\ge1}L_k\mathbf1_{\{S_H(k)\le h\}},\qquad
 V_h^{(w)}=\sum_{k\ge1}L_k
   \mathbf1_{\{S_H(k)\le h,\ S(\Gamma_k)\le w\}}.
\]
Thus $V_h^{(w)}$ counts a block only when its entire past, through that
block's endpoint, remains in the lateral interval $[-w,w]$. Later blocks
do not change whether an earlier block is counted.

\begin{corollary}[Consequences under the weaker finite inputs]
\label{R:alt:pv:laws}
Let $\gamma_0$ be the source port, $\gamma_n$ the $n$th visited vertex
of the infinite concatenation, and $Y_n$ its forward height from the source.
Almost surely,
\[
 Y_n\ge n^{3/4-o(1)},\qquad
 |\gamma_n-\gamma_0|\le n^{17/18+o(1)}.
\]
For the normalized finite bridge law summed over heights with weight
$\rho^Le^{-L/n}$, its height is at least $n^{3/4-\varepsilon}$ with
probability tending to one, for every $\varepsilon>0$.
The unnormalized mass of half-plane arches of length greater than $n$
is at most $n^{-3/16+o(1)}$.
The completed-vertex rewards just defined satisfy, for some fixed $C>0$,
\begin{equation}\label{R:alt:pv:prefix}
 \mathbb E V_h=h^{4/3+o(1)},\qquad
 \mathbb E V_h^{(C h\log^2h)}=h^{4/3+o(1)}.
\end{equation}
Consequently the expected number of vertices before first exit from a disk
of radius $R$ is at least $R^{4/3-o(1)}$.
\end{corollary}
\begin{proof}
We give the estimates that distinguish these conclusions from the sharp
spatial law.  At a dyadic vertex time $n$, let
$m=\lfloor n^{9/16-\varepsilon}\rfloor$.
The length tail in \eqref{R:alt:pv:tails}, a union bound above $n$, and
Markov's inequality for lengths truncated at $n$ imply
\[
 p\Bigl(\sum_{i\le m}L_i>n\Bigr)\le n^{-\varepsilon+o(1)}.
\]
The probability that all of those $m$ heights are below
$n^{3/4-2\varepsilon/\beta}$ is stretched exponentially small, by the
matching height tail.  Once these pieces have been completed, every later
vertex is above their slabs.  Borel--Cantelli on dyadics, followed by
monotonicity between dyadics, proves the asserted height lower bound.

Conversely take $m=\lceil n^{17/24+\varepsilon}\rceil$.
The last estimate in \eqref{R:alt:pv:tails} and exponential Markov give
\[
 p\Bigl(\sum_{i\le m}L_i\le n\Bigr)
 \le e\bigl(\mathbb E_p e^{-L/n}\bigr)^m
 \le \exp(-n^{\varepsilon/2})
\]
for large $n$.  The diameter-tail bound, a union bound and truncated first
moments show that $\sum_{i\le m}D_i\le m^{1/\beta+\varepsilon}$ except
with a power-decaying probability.  These bounds are summable on dyadics,
and the union of the first $m$ blocks contains every vertex up to time $n$.
Sending $\varepsilon$ to zero gives $(17/24)/\beta=17/18$.

Including the empty bridge, the height-summed tilted law has normalizer
\[
 Z_n=(1-\mathbb E_pe^{-L/n})^{-1}\ge n^{9/16-o(1)}.
\]
The bound uses the upper length tail.  Excluding the empty bridge replaces $Z_n$
by $Z_n-1$ and does not change the conclusion.  Its part of height at most $h$ is bounded by
$\sum_{j\le h}b_j=h^{\beta+o(1)}$.
Take $h=n^{3/4-\varepsilon}$.  For arches, split at height
$h=n^{3/4}$: the part leaving that strip has mass $O(b_h)$;
the remaining long arches have mass at most
$n^{-1}h^{13/12+o(1)}$.  Both contributions have exponent $-3/16$.

For the completed prefix, put $j_h=\mathbb E_p[L;H=h]$ and mark the
start height $l$ of each counted block.  The expected count is exactly
\[
 \mathbb E V_h=\sum_{l+j\le h}b_l j_j.
\]
The cumulative bounds in Proposition~\ref{R:alt:pv:reward} prove the upper
bound.  For the lower restrict $l\le h/2$ and apply the retained reward
at scale $r=h^{1-\delta}$, so eventually $j\le h/2$.
The prefix masses with $l\le h/2$ total $h^{\beta+o(1)}$.
Their contribution with lateral excursion greater than $h\log^2h$ is
superpolynomially small by \eqref{R:alt:pv:lateral}; the retained last
block has excursion at most $r\log^4r=o(h)$.
For each marked block, its preceding prefix stays within $h\log^2h$ of
the source and the marked block stays within $o(h)$ of its own start.
Hence the whole prefix through that block is confined to
$[-C h\log^2h,C h\log^2h]$ for one fixed $C$. These are precisely
contributions to $V_h^{(C h\log^2h)}$, of expected total mass at least
$h^{\beta+o(1)}r^{d-\beta-o(1)}$. Let $\delta\downarrow0$ and use
$V_h^{(w)}\le V_h$ to obtain both parts of \eqref{R:alt:pv:prefix}.
Taking $h$ a sufficiently small multiple of $R/\log^3R$ puts every retained
prefix inside the disk. All its counted vertices occur before first exit,
which proves the last assertion.
\end{proof}

\subsection{Completed-prefix mass in a fixed-proportion ball}
\label{R:alt:pd}

The disk transfer representation supplies a stronger confinement input
\cite[Theorem~1.1]{RenewalDiskCompanion}.  In its unit-edge embedding a layer has height
$3/2$, and its height-$h$ bridge masses satisfy
\begin{equation}\label{R:alt:pd:input}
 b_h\asymp(1+h)^{-1/4},\qquad m_h=h^{13/12+o(1)}.
\end{equation}
For every sufficiently large fixed $C$, the lower bound for $m_h$ remains
valid when the entire bridge has lateral wandering at most $Ch$.
Multiplying the preceding triangular-side-one coordinates by $\sqrt3$
gives this unit-edge embedding: the layer distance becomes $3/2$.
Vertex counts and critical weights are unchanged.  We keep layer height integral
in each embedding.

\begin{proposition}[A fixed-ball completed-prefix expectation]
\label{R:alt:pd:prefix}
Under the irreducible law from \eqref{R:alt:pd:input}, let
$\mathcal V_R$ be the vertices through the last block ending at cumulative
layer height at most $R$, excluding the overshooting block, and let
$\mathcal V_\infty$ be the full infinite trace.  With the starting port as
origin,
\[
 \mathbb E|\mathcal V_R|=R^{4/3+o(1)},\qquad
 \mathbb E|\mathcal V_R\cap B(0,R)|=R^{4/3+o(1)},
\]
and almost surely
\[
 |\mathcal V_\infty\cap B(0,R)|\le R^{4/3+o(1)}.
\]
Changing the ball radius by any fixed positive factor preserves the
expectation exponent.
\end{proposition}
\begin{proof}
Again put $j_h=\mathbb E_p[L;H=h]$ and $\beta=3/4$.
The renewal and reward identities with \eqref{R:alt:pd:input} give
\[
 B(1/R)\asymp R^\beta,\quad
 \sum_hm_he^{-h/R}=R^{25/12+o(1)},\quad
 \sum_hj_he^{-h/R}=R^{7/12+o(1)}.
\]
They imply $\sum_{h\le R}j_h=R^{7/12+o(1)}$.
For the lower bound, evaluate at $R^{1-\varepsilon}$: the tail beyond $R$
is exponentially negligible because $j_h\le m_h$ has polynomial growth;
then let $\varepsilon\downarrow0$.
Consequently $\sum_{l+h\le R}b_lj_h=R^{4/3+o(1)}$, proving the first
expectation and the upper bound for the second.

The fixed-ball lower bound needs the confined bridge estimate, rather than
only these transforms.  Choose $c_0>0$ so small that each height-$m$ bridge
with $c_0R\le m\le2c_0R$ and lateral wandering at most $Cm$ lies in
$B(0,R)$ and has height at most $R$.  The sum of their length masses is
\begin{equation}\label{R:alt:pd:marked}
 \sum_{c_0R\le m\le2c_0R}
   \sum_{\substack{\gamma:H(\gamma)=m\\\gamma\subset B(0,R)}}
       L(\gamma)\rho^{L(\gamma)}
 \ge R^{25/12-o(1)}.
\end{equation}
Under the iid irreducible law, the left side counts visits in prefixes
ending at such renewal heights, with a visit repeated at later permissible
prefix endings.  For a vertex in a block ending at height $s\le R$, condition
on that block and its past.  Dropping every future confinement condition,
the expected number of repetitions at later renewal endings is at most
$\sum_{h\le R}b_h=O(R^\beta)$, by independence of the future pieces.
Every counted vertex belongs to $\mathcal V_R\cap B(0,R)$.
Tonelli's theorem therefore bounds the left side of
\eqref{R:alt:pd:marked} by
$CR^\beta\mathbb E|\mathcal V_R\cap B(0,R)|$.
Since $25/12-\beta=4/3$, this proves the lower bound.  The same choice of a
smaller $c_0$ handles any fixed positive multiple of the ball radius.

For the almost-sure statement, the bounded-factor height deficit gives
$p(H>R)\asymp R^{-\beta}$: its upper bound is immediate, and after using
that upper bound the contribution of $H\le\eta R$ to the deficit at $R$
is at most $C\eta^{1-\beta}R^{-\beta}$.  Choose a fixed small $\eta$;
the remaining contribution gives $p(H>\eta R)\ge cR^{-\beta}$.
Replacing $R$ by $R/\eta$ proves the lower bound at $R$.
Splitting according to height and using the cumulative reward gives
$p(L>t)\le t^{-9/16+o(1)}$.
The first $k=\lceil R^{\beta+\varepsilon}\rceil$ blocks reach a height
larger than the greatest layer meeting $B(0,R)$ except with probability
$\exp(-cR^\varepsilon)$.  Their total length is at most
$R^{4/3+3\varepsilon(4/3)/\beta}$ except with probability
$R^{-2\varepsilon+o(1)}$, by a union bound for a length above that threshold
and Markov's inequality for truncated lengths.  After those blocks, the
path remains above the ball.  The estimates are summable on dyadic radii;
monotonicity and $\varepsilon\downarrow0$ complete the proof.
\end{proof}
